\subsection{The three attachment conditions}

\begin{definition}[Known: the manuscript's variance]\label{F:definitions}
Let $i:U\hookrightarrow X$ be a full sieve of small categories, let
$T=X\setminus U$, and let $\overline U$ be the full cosieve generated by
$U$. For $t\in T$ put
\[
 H_t:U^{\op}\longrightarrow\Set,\qquad H_t(u)=\Hom_X(u,t).
\]
The immersion is $W$-dwy\'erienne if every nonempty $H_t$ is a totally
$W$-aspherical presheaf. It is totally $W$-dwy\'erienne if every nonempty
$U/t$ is a totally $W$-aspherical category. These are precisely the
conditions that $U\to\overline U$ be strictly, respectively totally,
$W$-aspherical in XII ms.~109. A classical Dwyer map additionally requires
that $U\to\overline U$ have a right adjoint. A $W$-cofibration is a functor
$i:U\to X$ such that
\[
 i_!:U\backslash\Cat\longrightarrow X\backslash\Cat,
 \qquad (U\to C)\longmapsto(X\to X\amalg_U C),
\]
preserves morphisms whose underlying functors belong to $W$.
\end{definition}

The last definition is exactly that of an $h$-cofibration, with weak
equivalences $W$, in \citet[Definition~1.1]{F_BataninBerger2017}. We use
``classical'' to retain Thomason's right-adjoint requirement
\citep[Definition~4.1]{F_Thomason1980}; Cisinski's pseudo-Dwyer maps have
only a retraction with a relative natural transformation
\citep[Definition~1]{F_Cisinski1999}. The two requirements must not be
identified.

\begin{proposition}[Proved: cobase change and the forward part of
Th\'eor\`eme~3]\label{F:pushout}
Let $W$ be a fundamental localizer. Both classes of dwy\'erienne sieves
in Definition~\ref{F:definitions} are stable under arbitrary cobase change
in $\Cat$. Every $W$-dwy\'erienne sieve is a $W$-cofibration. In a pushout
square
\[
\begin{tikzcd}
 U\arrow[r,"i",hook]\arrow[d,"p"']&X\arrow[d,"q"]\\
 C\arrow[r,"i'",hook]&X\amalg_U C,
\end{tikzcd}
\]
the implications $p\in W\Rightarrow q\in W$ and
$i\in W\Rightarrow i'\in W$ hold. The assertions apply to arbitrary
small $C$ and arbitrary $p$, with no injectivity or fullness assumption
on $p$.
\end{proposition}

\begin{proof}
The pushout has old part $C$, complementary full subcategory $T$, and
crossing presheaves $p_!H_t$. For another functor $r:C\to D$, the
canonical isomorphism
$(rp)_!H_t\cong r_!p_!H_t$ proves preservation of totality of the
presheaf. If $U/H_t$ is totally aspherical, the canonical functor
\[
 U/H_t\longrightarrow C/(p_!H_t)
\]
is ordinary final, so its target is again totally aspherical by the
final-image criterion. An empty weight stays empty. A nonempty weight
stays nonempty, since any element gives a representative in the value
of its Kan extension at its image object. Thus the new boundary is the
same full subcategory $\dot T=\overline U\cap T$ of $T$.

For completeness, we give the gluing argument needed for the
cofibration assertion. The two cosieves $\overline U,T\subset X$ cover
$X$, and their intersection is $\dot T$. Their nerves form a literal
union. The same is true after any cobase change. The map
$U\to\overline U$ is $W$-aspherical: its slices over objects of $U$
have terminal objects, and the other slices are the categories of
elements of the nonempty $H_t$. The relative Theorem~A axiom implies
$U\to\overline U\in W$.

We recall explicitly why the forward gluing assertion used here is
valid for every fundamental localizer. Replace a union of two cosieves
$V_1\cup_{V_0}V_2$ by the Grothendieck construction of the covariant
span $V_1\leftarrow V_0\to V_2$. The functor from that construction to
the union is a weak homotopy equivalence. Over an object belonging to
just one of the cosieves its comma category has the corresponding
terminal object. Over an object of the intersection its comma
category is the Grothendieck construction of the three slice
categories, each with its compatible terminal object; contraction to
the span of these three terminal objects contracts its nerve.
An objectwise $W$-equivalence of such spans induces a $W$-equivalence
on their Grothendieck constructions, by the relative Theorem~A axiom
over the indexing category. All comparison functors to the unions
belong to $W_\infty\subseteq W$. Consequently a morphism of the three
pieces which belongs to $W$ piece by piece induces a functor in $W$
on the unions. Only this forward gluing statement is used.

Now take a morphism $f:C\to D$ under $U$ in $W$. Write $\overline C$
and $\overline D$ for the generated cosieves in the two pushouts.
The inclusions $C\to\overline C$ and $D\to\overline D$ belong to
$W$ by the first paragraph and Theorem~A. Thus
$\overline C\to\overline D$ belongs to $W$ by two-out-of-three.
Glue this map with the identities on $\dot T$ and $T$. The preceding
argument proves
$X\amalg_U C\to X\amalg_U D\in W$, which is the full
$W$-cofibration assertion, including every preliminary change of base.
The final implication $i\in W\Rightarrow i'\in W$ follows from
Proposition~\ref{F:couniversal} below.
\end{proof}

\begin{proposition}[Proved: couniversality in the fundamental-localizer
setting]\label{F:couniversal}
For every fundamental localizer $W$, a functor is a couniversal
$W$-equivalence if and only if it is both a $W$-cofibration and a member
of $W$. The class of $W$-cofibrations is closed under composition,
cobase change, and retracts in the arrow category.
\end{proposition}

\begin{proof}
The closure assertions follow by composing the pushout functors, by
the associativity isomorphism for pushouts, and by applying a functor
to a retract diagram, respectively. If every cobase change of $i$
belongs to $W$, the natural square comparing two such changes gives
the preservation property by two-out-of-three; the identity change
also gives $i\in W$.

Conversely let $i:U\to X$ be a $W$-cofibration in $W$, and let
$p:U\to C$ be arbitrary. Factor $p=w j$, with $j:U\to M$ a
Thomason cofibration and $w:M\to C$ in $W_\infty$, using
\citet[Theorem~4.9]{F_Thomason1980}. Every Thomason cofibration is a
retract of a classical Dwyer map by
\citet[Proposition~2]{F_Cisinski1999}. Proposition~\ref{F:pushout}
and retract closure therefore make $j$ a $W$-cofibration. Applying
its pushout functor to $i\in W$ gives
$M\to X\amalg_U M$ in $W$. Applying the pushout functor of $i$
to $w\in W_\infty\subseteq W$ gives
$X\amalg_U M\to X\amalg_U C$ in $W$. Since $w$ is also in
$W$, the commuting square and two-out-of-three imply
$C\to X\amalg_U C\in W$.

This proof uses the fundamental-localizer hypothesis. It does not
turn the same conclusion into a theorem for an arbitrary class of
functors satisfying only two-out-of-three. For $W_\infty$ the result
also follows directly from the left proper case of
\citet[Lemma~1.10]{F_BataninBerger2017}.
\end{proof}

\begin{proposition}[Proved: the conditional equivalences and the
boundary corollary]\label{F:conditional-theorem3}
In the square of Proposition~\ref{F:pushout}, suppose that $i$ is
$W$-dwy\'erienne. Under the full axiom W(7 bis) of XII
ms.~76, the following conditions are equivalent:
\[
 i\in W,\qquad i'\in W,\qquad
 (\dot T\hookrightarrow T)\in W,\qquad
 i\text{ is a couniversal }W\text{-equivalence}.
\]
Under the full axiom W(7) of XII ms.~77 one also has
$p\in W\Longleftrightarrow q\in W$.
The equivalence between the first and last conditions requires
neither additional axiom.
\end{proposition}

\begin{proof}
Let $\overline C$ be the cosieve generated by $C$ in
$X\amalg_U C$. The boundary $\dot T$ is unchanged, and
$U\to\overline U$ and $C\to\overline C$ belong to $W$, by
Proposition~\ref{F:pushout}. Apply the cosieve-cover form of
W(7 bis) to each of
\[
 X=\overline U\cup_{\dot T}T,
 \qquad
 X\amalg_U C=\overline C\cup_{\dot T}T.
\]
It gives
$\overline U\to X\in W\Longleftrightarrow
\dot T\to T\in W\Longleftrightarrow
\overline C\to X\amalg_U C\in W$.
Two-out-of-three gives the first three stated equivalences.
The first and last are equivalent by
Propositions~\ref{F:pushout} and~\ref{F:couniversal}.

For the second assertion, the map between these two cover
squares is the identity on $\dot T$ and $T$. W(7), in its
cosieve-cover form, therefore identifies
$q\in W$ with $\overline U\to\overline C\in W$.
The commuting square whose other two sides are
$U\to\overline U$ and $C\to\overline C$ identifies the
latter condition with $p\in W$. This proves exactly
Th\'eor\`eme~3(a,b) and Corollaire~2 under their printed
extra hypotheses; Corollaire~1 and cobase-change
Corollaire~3 are strengthened by Proposition~\ref{F:pushout}.
\end{proof}

The reverse detection statements are not consequences of the
fundamental-localizer axioms. The acyclic noncontractible example
in Part~A explains why the full W(7 bis) cannot be assumed for
$W_\infty$. There are also typographical errors at
ms.~110 and 113: the pushout is $U'\amalg_U X$, with $X$, not $X'$,
as its original second entry.

\subsection{Composition of the total classes}

\begin{lemma}[Proved: operations on totally aspherical categories]
\label{F:total-operations}
For $W=W_\infty$, or $W=W_R$ for a nonzero commutative unital ring
$R$, connected components of discrete opfibrations over
a totally $W$-aspherical category are totally $W$-aspherical. Every
nonempty cosieve in such a category is totally $W$-aspherical. If
$B$ and every value of a small covariant functor $D:B\to\Cat$ are
totally $W$-aspherical, then $\int_BD$ is totally $W$-aspherical.
\end{lemma}

\begin{proof}
We use the action criterion proved in Part~C: a nonempty connected
category $A$ is totally $W_\infty$-aspherical precisely when each
connected component of every covariant category of elements over
$A$ is weakly contractible. To recall the essential point in its
necessity, the collage test first gives acyclicity. Every connected
covering of a category of elements is itself a discrete opfibration
over $A$. A covering corresponding to a nontrivial cyclic subgroup
of the fundamental group would have nonzero $H_1$, a contradiction.
Thus those components are simply connected and acyclic. For its
sufficiency, each comma category of a final functor out of $A$ is
such a connected category of elements, and Theorem~A applies.

Discrete opfibrations compose. A further covariant category of
elements over a component of one is therefore again a discrete
opfibration over $A$, after extending by the empty set on the other
components. The criterion proves the first assertion. A nonempty
cosieve $S\subset A$ is the category of elements of the covariant
indicator functor which is a singleton on $S$ and empty elsewhere.
Totality forces common successors, so $S$ is connected: two objects
of $S$ have a common successor in $A$, and that successor is still
in $S$. The second assertion follows.

Put $E=\int_BD$. Its nerve is nonempty and connected because the
base and all fibers are nonempty and connected. Let $H:E\to\Set$
be covariant. Its category of elements is canonically the
Grothendieck construction of the categories
$E_b=\operatorname{El}(H|_{D(b)})$, with the transition functors
induced by $D$ and $H$. Every connected component of $E_b$ is
contractible by the action criterion for $D(b)$. The natural maps
$E_b\to\pi_0 E_b$ are therefore weak homotopy equivalences.
Their Grothendieck constructions give a weak homotopy equivalence
\[
 \operatorname{El}(H)\longrightarrow
 \operatorname{El}\bigl(b\mapsto\pi_0 E_b\bigr).
\]
This is the usual objectwise invariance of Grothendieck
constructions, or the relative Theorem~A axiom over $B$.
Every component of the target is contractible, by the criterion for
$B$. Therefore every component of $\operatorname{El}(H)$ is
contractible. The criterion applied to $E$ proves the last assertion.
All the constructions use the actual transition functors. No
passage to an ind-object in a homotopy category is involved.
For $W_R$, use the homological action criterion of
Theorem~\ref{D:R-action}. The same argument replaces every
occurrence of ``contractible'' by ``$R$-acyclic'' and every
weak equivalence of Grothendieck constructions by an
$R$-homology equivalence. The common-successor condition
still holds by the same nonzero-$R$ collage obstruction.
This proves all three assertions in that case as well.
\end{proof}

\begin{theorem}[Proved: composition of total functors and total
immersions for weak and homological equivalences]\label{F:total-composition}
Totally $W$-aspherical functors and totally
$W$-dwy\'erienne sieves are closed under composition when
$W=W_\infty$ or $W=W_R$ for a nonzero commutative unital ring $R$.
\end{theorem}

\begin{proof}
Let $f:X\to Y$ and $g:Y\to Z$ be totally aspherical functors.
For $z\in Z$, consider the covariant diagram on $Y/z$
\[
 (y,\beta:g(y)\to z)\longmapsto X/y.
\]
An arrow $h:y\to y'$ in $Y/z$ acts by postcomposing
$\alpha:f(x)\to y$ with $h$. Its Grothendieck construction $E_z$
is totally aspherical by Lemma~\ref{F:total-operations}. There is
a functor in $\Cat$
\[
 r:E_z\longrightarrow X/z,
 \qquad (y,\beta,x,\alpha)\longmapsto
 (x,\beta\,g(\alpha)).
\]
Its section $s:X/z\to E_z$ sends $(x,\gamma)$ to
$(f(x),\gamma,x,\id_{f(x)})$. For $c=(x,\gamma)\in X/z$
and $d=(y,\beta,x',\alpha)\in E_z$, a morphism $s(c)\to d$
is uniquely determined by a morphism $k:x\to x'$ satisfying
$\beta g(\alpha f(k))=\gamma$; its $Y$-component is
$\alpha f(k)$. This is exactly a morphism $c\to r(d)$.
Consequently $s\dashv r$, with $rs=\id$, and $r$ is ordinary
final. Totality passes to its target $X/z$.

For sieves $U\subset X\subset Y$, fix $y\in Y$ with $U/y$
nonempty. Then $X/y$ is totally aspherical. In $X/y$ let $S_y$
be the full subcategory of $(x,\beta:x\to y)$ for which $U/x$
is nonempty. It is a nonempty cosieve, hence is totally aspherical
by Lemma~\ref{F:total-operations}. The diagram
$(x,\beta)\mapsto U/x$ on $S_y$ has totally aspherical values.
Its Grothendieck construction maps to $U/y$ by composition,
with the same adjoint section as above. Therefore $U/y$ is
totally aspherical. This proves exactly the condition required
for the composite sieve. Empty boundary slices impose no condition.
\end{proof}

\begin{theorem}[Proved: strict composition for every fundamental
localizer]\label{F:strict-composition}
If $f:X\to Y$ is strictly $W$-aspherical, restriction
$f^*:Y^{\wedge}\to X^{\wedge}$ takes totally
$W$-aspherical presheaves to totally $W$-aspherical
presheaves. Consequently strictly $W$-aspherical functors
are closed under composition for every fundamental
localizer $W$.
\end{theorem}

\begin{proof}
Let $H:Y^{\op}\to\Set$ be totally aspherical and fix
an arbitrary functor $p:X\to X'$. For $y\in Y$ set
\[
 Q_y=p_!\Hom_Y(f(-),y)\in (X')^{\wedge}.
\]
These presheaves form a covariant functor
$Q:Y\to(X')^{\wedge}$. Every $Q_y$ is totally
$W$-aspherical, by the hypothesis on $f$ and
transitivity of Kan extension. Form the small collage
$Z$ with full subcategories $X',Y$, no arrows from
$Y$ to $X'$, and
$\Hom_Z(x',y)=Q_y(x')$. Composition is given by the
contravariance in $x'$ and the covariance in $y$.
Write $j:Y\hookrightarrow Z$ for the complementary
cosieve and $F=j_!H$. Since $H$ is total, $Z/F$ is
$W$-aspherical. Its restriction to $Y$ is $H$, and
on $X'$ it is the presheaf
\[
 K=\int^{y\in Y} Q_y\times H(y)
   \cong p_! f^*H.
\]
The isomorphism uses co-Yoneda in $Y^{\wedge}$,
pointwise preservation of colimits by $f^*$, and
preservation of colimits by $p_!$. All are ordinary
presheaf colimits.

The inclusion $l:X'/K\hookrightarrow Z/F$ is a
full sieve. Its slice over an object $(x',k)$ of
$X'/K$ has terminal object $((x',k),\id)$.
Its slice over $(y,h)$, with $h\in H(y)$, is
canonically isomorphic to $X'/Q_y$: an element
$q\in Q_y(x')$ gives the unique source element
$[q,h]\in K(x')$ and the corresponding crossing
arrow $(x',[q,h])\to(y,h)$. Conversely each
crossing arrow has precisely these data. Morphisms
are the arrows of $X'$ carrying $q'$ to $q$;
the condition on $[q,h]$ then follows automatically.
Thus every slice of $l$ is $W$-aspherical. By the
relative Theorem~A axiom $l\in W$, and since
$Z/F\to e$ is in $W$, so is $X'/K\to e$.
The functor $p$ was arbitrary, proving totality of
$f^*H$.

If $g:Y\to Z_0$ is strictly aspherical, apply this
result to each totally aspherical presheaf
$\Hom_{Z_0}(g(-),z)$. Its restriction along $f$
is exactly $\Hom_{Z_0}(gf(-),z)$, proving the
composition assertion.
\end{proof}

\begin{theorem}[Proved: total-functor composition for every
fundamental localizer]\label{F:total-functor-general}
For every fundamental localizer $W$, totally $W$-aspherical
functors are closed under composition. For a small covariant
diagram $D:B\to\Cat$, if $B$ and every $D(b)$ are totally
$W$-aspherical, then $\int_BD$ is totally $W$-aspherical.
\end{theorem}

\begin{proof}
Let $f:X\to Y$ and $g:Y\to Z$ be totally
$W$-aspherical, and fix $z\in Z$. Define
\[
 f_z:X/z\longrightarrow Y/z,
 \qquad (x,\gamma:gf(x)\to z)\longmapsto(f(x),\gamma),
\]
where $X/z$ is formed using $gf$ and $Y/z$ using $g$.
For $(y,\beta:g(y)\to z)$, the comma category
$(X/z)/(y,\beta)$ is canonically isomorphic to $X/y$.
Indeed an object is $x$ and $h:f(x)\to y$, since the
remaining arrow is forced to be $\gamma=\beta g(h)$.
A morphism $k:x\to x'$ is subject exactly to
$h'f(k)=h$; its other required equation follows from
this one. Thus $f_z$ is totally $W$-aspherical.
Proposition~\ref{E:forward-subterminal} makes it strictly
$W$-aspherical. The terminal presheaf of $Y/z$ is total
because $g$ is total. By Theorem~\ref{F:strict-composition}
its restriction along $f_z$, the terminal presheaf of
$X/z$, is total. This proves totality of $gf$ for all
$z$, without an extra hypothesis on $W$.

For the second assertion put $E=\int_BD$ and let
$p:E\to B$ be its projection. For $b\in B$ define
\[
 r_b:E/b\to D(b),\qquad
 (b',x,\beta:b'\to b)\longmapsto D(\beta)(x).
\]
Its section $s_b:D(b)\to E/b$ sends $y$ to
$(b,y,\id_b)$. A morphism from $(b',x,\beta)$ to
$s_b(y)$ is exactly a morphism $D(\beta)(x)\to y$
in $D(b)$, so $r_b\dashv s_b$. Hence $s_b$ is
ordinary final. Since $D(b)$ is total, its final
image $E/b$ is total. Thus $p$ is a totally
$W$-aspherical functor. So is $B\to e$; their
composite is total by the first assertion. This
says precisely that $E$ is totally $W$-aspherical.
\end{proof}

\begin{theorem}[Proved: the exact fixed-localizer criterion for
total sieve composition]\label{F:total-sieve-cosieves}
For a fixed fundamental localizer $W$, totally
$W$-dwy\'erienne sieves are closed under composition
if and only if every nonempty full cosieve of every
totally $W$-aspherical category is totally
$W$-aspherical. For $W_n$ with $n\geq2$, composition
of totally $W_n$-dwy\'erienne sieves holds whenever
both inclusions are between finite categories
(finitely many arrows).
\end{theorem}

\begin{proof}
Assume the cosieve condition, and take totally
$W$-dwy\'erienne sieves $U\subset X\subset Y$.
For $y$ with $U/y$ nonempty, let $S_y$ be the
full subcategory of $X/y$ consisting of the
objects $(x,\beta)$ for which $U/x$ is nonempty.
It is a nonempty cosieve in the total category
$X/y$, so is total by hypothesis. Every value
of the covariant diagram $(x,\beta)\mapsto U/x$
on $S_y$ is total. Theorem~\ref{F:total-functor-general}
makes its Grothendieck construction total. The
composition functor from this construction to
$U/y$ has the adjoint section explicitly given
in the proof of Theorem~\ref{F:total-composition};
it is therefore final. Thus $U/y$ is total,
proving composition of the sieves.

Conversely let $A$ be totally $W$-aspherical and
let $V\subset A$ be a nonempty full cosieve.
Form $X$ with a full old copy of $A$, a full
lower copy $V_0$ of $V$, and crossing hom sets
\[
 X(v_0,a)=A(v,a),\qquad X(a,v_0)=\varnothing.
\]
All compositions are induced by composition in
$A$ and $V$. The sieve $V_0\hookrightarrow X$
is classical Dwyer: its generated cosieve is
$V\times[1]$, and projection to the lower
copy is right adjoint to its inclusion.
The functor $r:X\to A$ which forgets the
copy is left adjoint to the old inclusion
$j:A\to X$, because
$X(x,j(a))=A(r(x),a)$ naturally. Thus $j$
is final, and $X$ is total as a final image
of $A$. Adjoining a terminal object $t$ gives
a totally $W$-dwy\'erienne sieve
$X\hookrightarrow X^{\triangleright}$: its
only new boundary category is $X$. Under
the assumed composition property, the
composite from $V_0$ is totally dwy\'erienne.
Its nonempty boundary category over $t$
is exactly $V$. Therefore $V$ is total,
proving necessity.

For the finite assertion, every nonempty boundary
slice of either inclusion is a finite category
which is totally $W_n$-aspherical. By
Corollary~\ref{D:truncated-finite} it is filtered,
hence totally $W_\infty$-aspherical. Both
inclusions are therefore totally
$W_\infty$-dwy\'erienne. Their composite has
that property by Theorem~\ref{F:total-composition},
and consequently is totally $W_n$-dwy\'erienne.
\end{proof}

\begin{theorem}[Refuted: low-degree sieve composition]
\label{F:low-sieve-composition}
For each of $W_0$ and $W_1$, neither the
$W$-dwy\'erienne sieves nor the totally
$W$-dwy\'erienne sieves are closed under
composition. In each case there is a finite
counterexample whose first factor is classical
Dwyer, whose second factor is totally
$W$-dwy\'erienne, and whose final target has
six objects. Thus the assertion of total
sieve composition for every fundamental
localizer is false, whereas total functor
composition holds by
Theorem~\ref{F:total-functor-general}.
\end{theorem}

\begin{proof}
For $W_0$ take $V$ to be the discrete category
on $a,b$. For $W_1$ take the two-object
category with exactly two nonidentity arrows
$\alpha,\beta:a\to b$. Its nerve realization
is a circle: the two edges have the same
endpoints, and there are no nondegenerate
simplices in higher dimensions. In particular,
$V$ is not $W$-aspherical in the indicated
case; the second choice is connected but
has fundamental group $\mathbb Z$.

Put $B=V^{\triangleleft}$, with new initial
object $o$, and apply the lower-copy
construction in the preceding proof to
$V\subset B$. Explicitly, $X$ has the old
objects $o,a,b$, the lower objects $a_0,b_0$,
the arrows of $B$ and of $V_0$, and
\[
 X(v_0,b')=B(v,b'),\qquad
 X(b',v_0)=\varnothing
 \quad(v\in V,\ b'\in B).
\]
The only compositions not already internal
to the two copies are precomposition and
postcomposition of these crossing arrows
in $B$. Let $Y=X^{\triangleright}$, with
new terminal object $t$. These specifications
give every object, arrow, and composite.

The inclusion $a:V_0\hookrightarrow X$
is classical Dwyer by the explicit right
adjoint above. The retraction $r:X\to B$
has a natural transformation
$\id_X\Rightarrow jr$, identity on the
old copy of $B$ and equal to the crossing
identity $v_0\to v$ on the lower copy.
Thus $BX\simeq BB$, which is contractible
since $o$ is initial. Theorem~\ref{D:zero-one}
makes $X$ totally $W_0$- and totally
$W_1$-aspherical. Consequently
$b:X\hookrightarrow Y$ is totally
$W$-dwy\'erienne, and both factors are
also strictly $W$-dwy\'erienne.

The composite's boundary category at $t$
is $V_0/t\cong V$, since there is exactly
one arrow from each old object to $t$.
For total dwy\'erienness this category
already fails the identity test. For strict
dwy\'erienness the crossing presheaf on
$V_0$ is its terminal presheaf, whose
category of elements is again $V$; its
identity test fails as well. This proves
both failures of composition.

The program \texttt{code/f\_low\_truncation.py}
constructs these categories, checks every
composition and associativity equation,
checks the right adjoint and its triangle
identities, checks the stated natural
transformation, and constructs the boundary
slice with its isomorphism to $V$. The two
targets have respectively $6$ objects and
$15$ arrows, and $6$ objects and $21$ arrows.
Their normalized simplex counts are
$(6,9,4)$ and $(6,15,16,6)$, respectively.
Both targets have trivial reduced integral
homology; the composite boundary has
$H_0=\mathbb Z^2$ in the first example
and $(H_0,H_1)=(\mathbb Z,\mathbb Z)$
in the second. Complete target arrow and
composition tables are included in the
program's output.
\end{proof}

\begin{lemma}[Proved: nonempty cosieve restrictions for homology]
\label{F:cosieve-homology}
Let $R$ be a nonzero commutative unital ring, let
$H:U^{\op}\to\Set$ be totally $W_R$-aspherical,
and let $V\subset U$ be a full cosieve. If
$H|_V$ is nonempty, it is totally $W_R$-aspherical.
\end{lemma}

\begin{proof}
Fix $p:V\to C$ and set $K=p_!(H|_V)$.
Adjoin an initial object $0$ to $C$, obtaining
$B=C^{\triangleleft}$. There is a functor
$q:U\to B$ agreeing with $p$ on $V$ and
sending every object outside $V$ to $0$.
There are no arrows from $V$ to its complement,
and all remaining mixed arrows go to the unique
arrows from $0$, so this is well defined.
For $F=q_!H$ we have $F|_C=K$ and
$F(0)=\pi_0(U/H)=\{*\}$. The latter equality
follows by testing totality of $H$ on $U\to e$,
since a discrete set is $R$-acyclic only if it
is a singleton. The category $B/F$ is exactly
the initial cone on $E=C/K$.

Define a covariant functor $L:B\to\Set$ by
$L(0)=\{a,b\}$ and $L(c)=\{*\}$ for
$c\in C$. Every arrow from $0$ collapses
the two elements, and every arrow of $C$
acts identically. Adjoin an object $w$ with
outgoing hom functor $L$, obtaining the
collage $B_L$, and let $r:B\to B_L$ be
the inclusion. The category of elements of
$L$ pulled back to $B/F$ consists of two
initial cone points on $E$. Its classifying
space is $\Sigma BE$. Because $K$ is nonempty,
this category is connected, and the ordinary
Kan extension has $(r_!F)(w)=\{*\}$.
The collage formula for the new category of
elements therefore gives
\[
 B(B_L/r_!F)\simeq
 \operatorname{hpo}\bigl(*\leftarrow
     \Sigma BE\to B(B/F)\bigr)
 \simeq\Sigma^2 BE.
\]
Here $B(B/F)$ is contractible, and the last
display uses the ordinary unreduced
suspension. Since $F$ is totally $W_R$-
aspherical, the left side is $R$-acyclic.
The reduced suspension homology isomorphism
and nonemptiness of $E$ imply that $E$ is
nonempty, connected, and $R$-acyclic. The
functor $p$ was arbitrary. Thus $H|_V$
is totally $W_R$-aspherical.
\end{proof}

\begin{theorem}[Proved: composition of homological dwy\'erienne
immersions]\label{F:strict-sieve-homology}
For every nonzero commutative unital ring $R$, the
$W_R$-dwy\'erienne sieves are closed under composition.
\end{theorem}

\begin{proof}
Let $U\subset X\subset Y$ be two such sieves,
and let $y\in Y$ receive an arrow from $U$.
The weight $H=X(-,y)$ is total on $X$.
Let $V\subset X$ be the cosieve generated by $U$.
Its restriction $H|_V$ is nonempty and is total
by Lemma~\ref{F:cosieve-homology}. The functor
$U\to V$ is strictly $W_R$-aspherical by
definition. Theorem~\ref{F:strict-composition}
therefore shows that $H|_U=Y(-,y)|_U$ is
totally $W_R$-aspherical. These are exactly
the nonempty crossing weights of the composite.
\end{proof}

\subsection{Discretely flat inclusions and strict separations}

\begin{definition}[Known: discretely flat]\label{F:DFdefinition}
For a small category $A$, let $\operatorname{disc}A$ denote the
simplicial category with its hom sets regarded as discrete
simplicial sets. A functor $i:U\to X$ is discretely flat if
\[
 \operatorname{disc}X\amalg_{\operatorname{disc}U}(-)
\]
preserves Dwyer--Kan equivalences of simplicial categories under
$\operatorname{disc}U$. It is a discretely flat cofibration when
it is also injective on objects, and a discretely flat inclusion
when it is also faithful. These are
\citet[Definition~3.4]{F_HORR2024}. A sieve satisfies both extra
injectivity conditions.
\end{definition}

\begin{proposition}[Proved: an exact weight test for discrete flatness]
\label{F:weight-test}
A sieve $i:U\hookrightarrow X$ is discretely flat if and only if,
for each $t\in X\setminus U$, the functor
\[
 H_t\otimes_U-:\operatorname{Fun}(U,\mathbf{sSet})
       \longrightarrow\mathbf{sSet}
\]
preserves objectwise weak homotopy equivalences. The coend is formed
degree by degree in sets. Every totally $W_\infty$-dwy\'erienne
sieve is discretely flat. Every sieve with discrete source is
discretely flat.
\end{proposition}

\begin{proof}
In a simplicial pushout of the sieve along
$p:\operatorname{disc}U\to C$, the object set is
$\operatorname{Ob}C\amalg\operatorname{Ob}T$. The hom simplicial
sets internal to $C$ are unchanged, those internal to $T$ are the
original discrete hom sets, and there are no arrows from $T$ to
$C$. The remaining hom simplicial sets are
\[
 (X\amalg_U C)(c,t)=H_t\otimes_U C(c,p(-)).
\]
This follows degree by degree from the ordinary collage pushout
formula; composition is induced by precomposition in $C$ and
postcomposition in $T$. If all the displayed weight functors
preserve weak equivalences, a Dwyer--Kan equivalence $C\to D$
under $U$ induces weak equivalences on all these hom objects.
Essential surjectivity is inherited from $C\to D$, together with
the identity on the unchanged objects of $T$. This proves
sufficiency.

For necessity, given a diagram $K:U\to\mathbf{sSet}$ form the
simplicial collage $C_K$ obtained from $\operatorname{disc}U$ by
adjoining an object $x$, setting $C_K(x,u)=K(u)$, and having no
arrows into $x$ except its identity. An objectwise weak equivalence
$K\to L$ induces a Dwyer--Kan equivalence $C_K\to C_L$,
identity on objects. After pushout, the map on hom objects from
$x$ to $t$ is exactly $H_t\otimes_U K\to H_t\otimes_U L$.
Discrete flatness forces this to be a weak equivalence.

We spell out the sufficient totality assertion. If a category $A$
is totally aspherical, ordinary colimits
$\operatorname{Fun}(A,\mathbf{sSet})\to\mathbf{sSet}$ preserve
objectwise weak equivalences. Indeed the augmentation from the
simplicial replacement to the ordinary colimit is a weak
equivalence: in each simplicial degree it is the map from the
nerve of the category of elements of a set-valued diagram to its
set of components, and all these components are contractible by
the action criterion. Taking the diagonal of the resulting
bisimplicial sets preserves these weak equivalences. The ordinary
colimit is thus naturally weakly equivalent to the homotopy
colimit, which preserves objectwise weak equivalences.
For $A=U/H_t$ we have the canonical degreewise isomorphism
\[
 H_t\otimes_U K\cong\colim_{(u,h)\in U/H_t}K(u).
\]
This proves the assertion for totally dwy\'erienne sieves.
If $U$ is discrete, the weight functor is simply
$\coprod_{u\in U}H_t(u)\times K(u)$, which plainly preserves
weak equivalences, including empty weights and empty $U$.
\end{proof}

\begin{proposition}[Proved: the restricted-action criterion for
discrete flatness]\label{F:DF-actions}
Let $H:U^{\op}\to\Set$, put $E=U/H$, and let
$\pi:E\to U$ be the projection. The weighted colimit
$H\otimes_U-:\mathbf{sSet}^{U}\to\mathbf{sSet}$
preserves objectwise weak equivalences if and only if
every connected component of
\[
 \operatorname{El}_{E}(L\pi),\qquad L:U\to\Set,
\]
is weakly contractible. Among weights satisfying this
condition, $H$ is totally $W_\infty$-aspherical if and
only if $E$ is nonempty and connected. Consequently a
discretely flat sieve is $W_\infty$-dwy\'erienne if and
only if each of its nonempty boundary categories $U/t$
is connected.
\end{proposition}

\begin{proof}
For a covariant set-valued functor $L$, define a
covariant simplicial diagram $PL$ by
\[
 PL(u)_n=
 \coprod_{u_0\to\cdots\to u_n\to u}L(u_0).
\]
The face and degeneracy maps are those of the usual
simplicial replacement. Its augmentation to the
discrete simplicial set $L(u)$ sends an element along
the indicated composite. For fixed $u$, $PL(u)$ is
the nerve of the category of triples
$(v,l,\alpha:v\to u)$, with the evident commuting
triangles as arrows. Each component has the terminal
triple $(u,l_u,\id_u)$ for its common image
$l_u\in L(u)$. Thus $PL\to L$ is an objectwise
weak equivalence.

The co-Yoneda isomorphism identifies
\[
 (H\otimes_U PL)_n
   \cong
 \coprod_{u_0\to\cdots\to u_n}
       L(u_0)\times H(u_n)
   \cong
 N\operatorname{El}_E(L\pi)_n.
\]
For the last identification, the element of $H(u_n)$
determines all its predecessors by contravariance,
and the element of $L(u_0)$ determines all its
successors by covariance. These identifications
respect every face and degeneracy. Also
$H\otimes_U L\cong\pi_0\operatorname{El}_E(L\pi)$.
Hence homotopical preservation by the weight forces
the map from this nerve to its set of components to
be a weak equivalence, proving necessity.

Conversely, for $K:U\to\mathbf{sSet}$ consider the
bisimplicial two-sided bar construction
\[
 B_{q,n}(H,U,K)=
 \coprod_{u_0\to\cdots\to u_q}
       H(u_q)\times K(u_0)_n.
\]
For fixed $n$, its augmentation to the set
$(H\otimes_U K)_n$ is precisely the map just
described for $L=K_n$. The assumed condition makes
that augmentation a weak equivalence of simplicial
sets in the $q$ direction. The diagonal therefore
gives a natural weak equivalence
$\operatorname{diag}B(H,U,K)\to H\otimes_U K$.
For an objectwise weak equivalence $K\to K'$, each
fixed-$q$ map of bars is a coproduct of copies of
objectwise weak equivalences and is therefore a weak
equivalence. Its diagonal is a weak equivalence as
well. Two-out-of-three proves the required
preservation by $H\otimes_U-$.

Taking $L$ constant at a singleton shows that every
component of $E$ itself is contractible. Under
nonempty connectedness, the restricted-action
criterion of Proposition~\ref{E:restricted-actions}
therefore proves total $W_\infty$-asphericity of $H$.
The converse follows from the identity test for
totality. Apply this argument weight by weight and
use Proposition~\ref{F:weight-test} for the final
assertion.
\end{proof}

\begin{example}[Proved: strict does not imply total, even among
discretely flat sieves]\label{F:strict-not-total}
Let $U$ have one object and endomorphism monoid
$M=\{1,e\}$, with $e^2=e$. Let the presheaf
$H$ have elements $i,h_1,h_2$, with $H(e)$
constant at $i$, and form $X$ by adjoining one
object $t$ with $\Hom_X(-,t)|_U=H$ and no
arrows from $t$ to $U$. Then $U\hookrightarrow X$
is $W_\infty$-dwy\'erienne and discretely flat,
but is not totally $W_\infty$-dwy\'erienne.
\end{example}

\begin{proof}
The presheaf counterexample in Part~E proves
totality of $H$ and failure of totality of $U/H$.
We recall its construction to make the new claim
about discrete flatness independently checkable.
Let $R=U(-,*)$ and $Q=\{e\}\subset R$.
Multiplication by $e$ is an idempotent retraction
of $R$ onto $Q$, and
$H=R\amalg_QR$. After any Kan extension the
two copies of $R$ remain representable and
the two copies of $Q\to R$ remain split
monomorphisms. The category of elements of the
image of $Q$ is filtered, while those of the
two images of $R$ have terminal objects; their
union is weakly contractible. Thus $H$ is total.
The objects $(i,h_1,h_2)$ of $U/H$ have
arrows from $i$ to each $h_j$, and the two
distinct $h_j$ have no common successor.
Consequently $U/H$ is not total.

For a covariant simplicial diagram $L:U\to
\mathbf{sSet}$, put $V=L(*)$ and let $d=L(e)$.
Then
\[
 H\otimes_U L\cong V\amalg_{\operatorname{im}d}V.
\]
The image of a simplicial idempotent is a
retract, so its inclusion in $V$ is a
simplicial monomorphism. If $L\to L'$ is
an objectwise weak equivalence compatible
with the idempotents, the induced map on
their images is a retract of the underlying
weak equivalence in the arrow category.
The two pushout diagrams therefore have
weakly equivalent entries and cofibration
legs. Their pushouts are weakly equivalent.
Proposition~\ref{F:weight-test} gives discrete
flatness. This also proves that the first
inclusion between the total and strict
dwy\'erienne classes is strict.
\end{proof}

\begin{proposition}[Known, with explicit consequence: discrete
flatness implies homotopical flatness]\label{F:DFh}
Every discretely flat sieve is a $W$-cofibration for every
fundamental localizer $W$. In particular it is an $h$-cofibration
for $W_\infty$.
\end{proposition}

\begin{proof}
For such an inclusion, \citet[Theorem~1.5, in its generalized
form]{F_HORR2024} show that the canonical comparison
\[
 NX\amalg_{NU}NC\longrightarrow N(X\amalg_U C)
\]
is a categorical equivalence, hence a weak homotopy equivalence.
Since $NU\to NX$ is a monomorphism, the source computes the
homotopy pushout of the three nerves. Compare this description
for a morphism $C\to D$ under $U$ in $W$. The homotopy pushout
comparison is a $W$-equivalence by objectwise invariance of the
Grothendieck construction of a span, as used in
Proposition~\ref{F:pushout}. The two comparison maps belong to
$W_\infty\subseteq W$. Two-out-of-three gives the claim.
\end{proof}

\begin{theorem}[Refuted: C035, even for a couniversal equivalence]
\label{F:C035}
A sieve which is a couniversal $W_\infty$-equivalence need not be
$W_\infty$-dwy\'erienne. The same counterexample works for every
nonzero homology localizer and every $W_n$ with $n\geq0$.
\end{theorem}

\begin{proof}
Let $X$ have objects $a,t,s$, identity arrows, two arrows
$\alpha,\beta:a\to t$, an arrow $\gamma:t\to s$, and an
arrow $\delta:a\to s$. The only composites of two nonidentity
arrows are
\[
 \gamma\alpha=\gamma\beta=\delta.
\]
Put $U=\{a\}$. This is a full sieve and generates all of $X$
as a cosieve. Its crossing presheaf at $t$ has two elements, so
$U/t$ is a discrete two-object category. Already the identity
test shows that this presheaf is not totally aspherical for any
of the localizers in the statement.

Let $p:U\to C$ be any functor, with selected object $c=p(a)$,
and put $P=X\amalg_U C$. The objects are those of $C$ together
with $t,s$. The new hom sets, for $d\in C$, are
\[
 P(d,t)=C(d,c)\times\{\alpha,\beta\},\qquad
 P(d,s)=C(d,c),\qquad P(t,s)=\{\gamma\}.
\]
Composition with $\gamma$ forgets the second coordinate. There
are no arrows from $t$ or $s$ into $C$, and the other hom sets
are as stated in $X$ or $C$. Define $r:P\to C$ by the identity
on $C$, by $r(t)=r(s)=c$, by projecting these new hom sets to
$C(d,c)$, and by $r(\gamma)=\id_c$. Let $j:C\to P$ be the
inclusion. Define $F:P\to P$ to be the identity on $C$ and to
send both $t,s$ to $s$, sending each crossing arrow to its
composite to $s$. There are natural transformations
\[
 \id_P\Longrightarrow F\Longleftarrow jr
\]
whose components outside $C$ are, respectively,
$(\gamma,\id_s)$ and the arrows $c\to s$ corresponding to
$\id_c$. Both transformations are identities on $C$. Naturality
at the two arrows to $t$ is precisely the displayed relation
$\gamma\alpha=\gamma\beta$. Thus $Bj$ is a homotopy
equivalence with inverse $Br$, for every $p$. This proves
couniversality for $W_\infty$ and hence for every fundamental
localizer. The failure of dwy\'erienne totality remains visible
in degree zero for the stated localizers. By
Proposition~\ref{F:weight-test} the sieve is also discretely flat.

The program \texttt{code/f\_check.py} verifies the seven-arrow
category, all $21$ composable triples, and the two natural
transformations. The universal proof above, rather than a finite
sample of targets, supplies the required quantifier over $p$.
\end{proof}

\begin{corollary}[Refuted: one-point and complementary-change
cofibration recognition]\label{F:onepoint-cofibration}
The ordinary-cofibration assertion at XII ms.~111
(C037) is false: a $W$-cofibration sieve whose
complementary full category is isomorphic to $e$
(one object and only its identity) need not be
$W$-dwy\'erienne. The
complementary-change recognition assertion (C038)
is also false: even requiring every change of the
complement to give a $W$-cofibration does not force
dwy\'erienness. These failures hold for $W_\infty$,
for $W_R$ with any nonzero commutative unital ring
$R$, and for every $W_n$, $n\geq0$.
\end{corollary}

\begin{proof}
Let $U=\{a\}$, and let $X$ have the two objects
$a,t$ and exactly two nonidentity arrows
$\alpha,\beta:a\to t$. Its source is discrete,
so Proposition~\ref{F:weight-test} makes $U\to X$
discretely flat, and Proposition~\ref{F:DFh}
makes it a $W$-cofibration for every fundamental
localizer. Its only boundary presheaf has two
elements, so fails the identity test for total
asphericity for all localizers listed above.

For every small category $T_1$ and the unique
functor $T_1\to\{t\}$, changing the complement
produces the collage with source $U$ and two
crossing arrows to each object of $T_1$, with
each transition preserving the two labels.
Its source is still discrete. The same two
propositions make every such changed sieve a
$W$-cofibration, including the empty-complement
case. This proves the failure of C038. In
particular, retaining full W(7) does not repair
these assertions, since nonzero homology
localizers satisfy that axiom.

The two-object category is the parallel-arrow
category explicitly constructed and checked by
\texttt{code/f\_low\_truncation.py}. Its nerve
is a circle, so this particular inclusion is
not couniversal for $W_\infty$ or for nonzero
homology localizers. The following proposition
treats that narrower variant precisely.
\end{proof}

\begin{proposition}[Proved: the exact one-point couniversality
criterion]\label{F:onepoint-couniversal}
Let $H:U^{\op}\to\Set$ be any presheaf on a small
category, and let $X$ be obtained by adjoining one
object $t$, with $X(u,t)=H(u)$, no arrows from
$t$ to $U$, and only the identity endomorphism
of $t$. Then $U\hookrightarrow X$ is a
couniversal $W_\infty$-equivalence if and only
if $H$ is totally $W_{\mathbb Z}$-aspherical.
For a nonzero commutative unital ring $R$, it
is a couniversal $W_R$-equivalence if and only
if $H$ is totally $W_R$-aspherical.
\end{proposition}

\begin{proof}
After arbitrary cobase change $p:U\to C$, the
new weight is $F=p_!H$. Write $E=C/F$ and
$\pi:E\to C$. The one-point collage formula
of Lemma~\ref{A:audit-cylinder} identifies
the specified inclusion of classifying spaces
with the inclusion into the homotopy cofiber
\[
 BC\longrightarrow
 BC\amalg^h_{BE}*.
\]
If $E$ is empty, this adds a separate point
and is not a homology equivalence over a
nonzero ring. If $E$ is nonempty, its relative
homology is $\widetilde H_{k-1}(BE;R)$.
Consequently this inclusion is an $R$-homology
equivalence exactly when $E$ is $R$-acyclic.
Quantifying over all $p$ proves the assertion
for $W_R$ and the necessity of total
$W_{\mathbb Z}$-asphericity in the first
assertion, including nonemptiness.

We prove the remaining sufficiency. First let
$F$ be any totally $W_{\mathbb Z}$-aspherical
presheaf on $C$. Its element category $E$
is nonempty, connected, and acyclic. Moreover
the homomorphism
\[
 \varphi:\pi_1(BE,e)\longrightarrow
 G=\pi_1(BC,\pi(e))
\]
induced by its projection is trivial. To
verify this, take $g\in\operatorname{im}\varphi$
and $K=\langle g\rangle\subset G$. On the
component of $C$ containing $\pi(e)$ take
the covariant set-valued local system
associated to the transitive $G$-set with
stabilizer $K$, and take the empty set on
the other components. Its pullback to $E$
is a covering action. The component through
the distinguished coset has fundamental
group $\varphi^{-1}(K)$. It is acyclic by
Proposition~\ref{E:restricted-actions}.
But $\varphi^{-1}(K)\to K$ is surjective,
since $K\subset\operatorname{im}\varphi$.
As $K$ is abelian, this surjection factors
through the zero first homology group of
that covering component. Hence $K=1$.
This applies to every $g$, proving the claim.

Consider the connected component of $BC$
to which $BE$ maps; all other components
are unchanged by the attachment. Van Kampen
and $\varphi=0$ show that adjoining the
cone on $BE$ preserves its fundamental
group. On universal covers the same
attachment adjoins a disjoint family of
cones on copies of $BE$, indexed by $G$:
the pullback of the universal cover to
$BE$ is trivial because $\varphi=0$.
The relative integral homology of the
map of universal covers is thus
\[
 \bigoplus_{G}\widetilde H_{k-1}(BE;\mathbb Z)=0.
\]
The map is a homology equivalence between
simply connected CW complexes, hence a
weak homotopy equivalence by the homological
Whitehead theorem. Together with the
fundamental-group isomorphism this proves
that the original cofiber inclusion is a
weak homotopy equivalence.

If $H$ is total $W_{\mathbb Z}$, each
$F=p_!H$ is total $W_{\mathbb Z}$, since
$q_!p_!H\cong(qp)_!H$ for every subsequent
functor $q$ with source $C$. The
preceding argument applies after every
cobase change, proving couniversality
for $W_\infty$.
\end{proof}

\begin{theorem}[Proved: a homotopical cofibration need not be
discretely flat]\label{F:h-not-DF}
There is a four-object full sieve $U\hookrightarrow X$ which
is a couniversal $W_\infty$-equivalence but is not discretely
flat. Consequently the inclusion of discretely flat sieves in
the class of $W_\infty$-cofibration sieves is strict.
\end{theorem}

\begin{proof}
Let $U=[1]=(0\xrightarrow{a}1)$. Adjoin objects $t,s$ to form
$X$. The new arrows to $t$ are
$\beta,\gamma:1\to t$ and $c:0\to t$, with
$\beta a=\gamma a=c$. Adjoin a terminal object $s$ to this
three-object category. These rules specify all arrows and all
compositions. The inclusion of $U$ is a sieve. The retraction
$X\to U$ fixes $U$ and sends $t,s$ to $1$. The endofunctor
fixing $U$ and sending $t,s$ to $s$ again gives
$\id_X\Rightarrow F\Leftarrow ir$, relative to $U$.
Products with $[1]$ preserve pushouts in $\Cat$, so a natural
transformation relative to $U$ passes to every pushout. Therefore
every pushout inclusion is a categorical homotopy equivalence.

To disprove discrete flatness, take the simplicial circle
$K=\Delta[1]/\partial\Delta[1]$ and its simplicial cone $CK$.
Define diagrams $L,M:U\to\mathbf{sSet}$ by
\[
 L(0)=K,\quad L(1)=CK,\quad L(a):K\hookrightarrow CK;
 \qquad
 M(0)=K,\quad M(1)=*,\quad M(a):K\to *.
\]
The natural transformation $L\to M$ is the identity at $0$
and the contraction $CK\to *$ at $1$, so is an objectwise
weak equivalence. The weight $H_t$ has one element at $0$,
two at $1$, and its two elements restrict to the same element.
Its coends are
\[
 H_t\otimes_U L=CK\amalg_K CK,\qquad
 H_t\otimes_U M=*\amalg_K *=*.
\]
The first realizes as the suspension of a circle, hence has
$H_2\cong\mathbb Z$; the second is a point. Thus the map is
not a weak equivalence. Proposition~\ref{F:weight-test}, or its
explicit simplicial-collage construction, disproves discrete
flatness. The finite category and its relative transformations
are checked in the accompanying program.
\end{proof}

\begin{proposition}[Proved: a strict extension of classical Dwyer
maps]\label{F:classical-strict}
The inclusion $\mathbb N\hookrightarrow\mathbb N^{\triangleright}$,
where $\mathbb N$ has its usual order and the target adds a new
terminal object, is totally $W$-dwy\'erienne for every
fundamental localizer $W$, but is not a classical Dwyer map.
\end{proposition}

\begin{proof}
The only new crossing category is $\mathbb N$, which is
filtered and therefore totally $W$-aspherical. A right adjoint
would assign to the new terminal object a terminal object of
$\mathbb N$, which does not exist. The generated cosieve is
the entire target, so no smaller neighborhood can avoid this
obstruction.
\end{proof}

Thus the proved inclusions, among sieves, are
\[
 \{\text{classical Dwyer}\}
 \varsubsetneqq
 \{\text{totally }W_\infty\text{-dwy\'erienne}\}
 \varsubsetneqq
 \{\text{discretely flat}\}
 \varsubsetneqq
 \{W_\infty\text{-cofibrations}\},
\]
and also
\[
 \{\text{totally }W_\infty\text{-dwy\'erienne}\}
 \varsubsetneqq
 \{W_\infty\text{-dwy\'erienne}\}
 \varsubsetneqq
 \{W_\infty\text{-cofibrations}\}.
\]
The first inclusion into the discretely flat class is strict as
well: Theorem~\ref{F:C035} belongs to that class and fails even
the strictly dwy\'erienne condition. In particular discretely
flat sieves are not contained in the strictly dwy\'erienne
class.

\subsection{The endpoint product is not a homotopical cofibration}

\begin{theorem}[Refuted: the proposed endpoint stability, in its
smallest standard test]\label{F:interval-failure}
Let $n\geq2$, let
\[
 I=\Phi_n=(0\to1\leftarrow2\to\cdots\leftarrow2n),
 \qquad j:\{0,2n\}\hookrightarrow I,
\]
and let $i:\{0\}\hookrightarrow[1]$ select the initial object.
The sieve
\[
 \alpha=i\mathbin{\square}j:
 A=(\{0\}\times I)\cup([1]\times\{0,2n\})
       \hookrightarrow P=[1]\times I
\]
is not a $W_\infty$-cofibration. It is not a $W$-cofibration
for any fundamental localizer $W$ in which a category of
homotopy type $S^2$ is not aspherical.
\end{theorem}

\begin{proof}
Write $l_k=(0,k)$ and $u_k=(1,k)$. The objects of $A$ are
$l_0,\ldots,l_{2n},u_0,u_{2n}$. Its nonidentity arrows form
an unoriented tree, so $BA$ is contractible. Also $BP$ is
contractible, being the product of two contractible nerves.

Define $C$ to have objects $x,z,a,b$, two arrows
$\lambda,\mu:x\to z$, arrows $r:z\to a$ and $v:z\to b$,
and their composites, with relations
\[
 r\lambda=r\mu,\qquad v\lambda=v\mu.
\]
There are no other nonidentity arrows. The full category on
$x,z$ has classifying space $S^1$; the full subcategories
obtained by adjoining $a$ and $b$ are its two terminal cones.
Their nerves cover $NC$ and intersect in that circle. Hence
$BC\simeq S^2$. Let $D=C^{\triangleright}$, with new terminal
object $s$, so $BD$ is contractible.

There is a functor $p:A\to D$ defined by
\[
 p(l_0)=z,\quad p(u_0)=b,\quad
 p(l_k)=a\ (1\leq k\leq2n),\quad p(u_{2n})=a.
\]
It sends $l_0\to l_1$ to $r$, $l_0\to u_0$ to $v$, and
every remaining nonidentity arrow to $\id_a$. This types
correctly because every remaining arrow has both endpoints
among the vertices sent to $a$. It is a weak equivalence,
since its source and target are contractible.

Compute the ordinary pushout $Q=P\amalg_A D$. Its complement
to $D$ consists of $u_1,\ldots,u_{2n-1}$, with the upper
alternating-interval arrows. The weight for $u_1$ before
pushout is subterminal on the full four-object subcategory
\[
 S=\{l_0,l_1,l_2,u_0\},\qquad
 u_0\leftarrow l_0\to l_1\leftarrow l_2.
\]
The functor $S\to\Psi=(z\to a,z\to b)$ induced by $p$
is ordinary final: its outgoing comma categories at $z,a,b$
are respectively the connected tree $S$, the arrow
$l_2\to l_1$, and the singleton $u_0$. The inclusion
$\Psi\to C$ is final. At $x$ its outgoing comma category
is the four-edge square with vertices
$(z,\lambda),(z,\mu),(a,r\lambda),(b,v\lambda)$;
at the other objects it has an initial object. Consequently
the pushed-forward crossing weight at $u_1$ is the
subterminal presheaf on $D$ supported on $C$.

For every $k\geq2$, all vertices of $A/u_k$ are sent to
$a$, all their arrows to $\id_a$, and $A/u_k$ is a nonempty
connected tree (or singleton). The pushed-forward weight is
therefore the representable $D(-,a)$. All transitions between
these latter weights are identities; the transition from
$u_2$ to $u_1$ is the canonical map $D(-,a)\to1_C$.
These statements can also be read directly from the coend
\[
 \Hom_Q(d,u_k)=\int^{v\in A}
       D(d,pv)\times\Hom_P(v,u_k).
\]
They describe every hom set and every composition of $Q$.

Let $R\subset Q$ be the full subcategory on $D$ and $u_1$.
It is exactly $C$ with two terminal cone points $s,u_1$,
with no arrow between the two points. Define $r:Q\to R$
to fix $R$ and to send each $u_k$, $k\geq2$, to $a$.
On the representable crossing hom sets it is the Yoneda
identification with $D(-,a)$, and on arrows among those
$u_k$ it is $\id_a$. The arrow $u_2\to u_1$ goes to the
unique arrow $a\to u_1$. For the full inclusion $e:R\to Q$
we have $re=\id_R$ and a natural transformation
$er\Rightarrow\id_Q$, identity on $R$, whose component
at $u_k$ is the arrow $a\to u_k$ represented by $\id_a$.
The above description of the transitions checks naturality
on every crossing and upper-interval arrow. Hence $Be$ is
a homotopy equivalence. The literal nerve union of the two
cones gives
\[
 BQ\simeq BR\simeq\Sigma BC\simeq S^3.
\]

If $\alpha$ were a $W$-cofibration, applying its pushout
functor to the morphism $p:A\to D$ under $A$, with source
the identity of $A$, would force $P\to Q$ to belong to
$W$. Since $P$ is contractible and $Q$ has type $S^3$,
this disproves the cofibration assertion whenever $S^3$
is not $W$-aspherical.

To obtain the stated stronger obstruction with $S^2$,
use instead the same assignment $p_0:A\to C$, without
adjoining $s$. The identical coend calculation shows
that $Q_0=P\amalg_A C$ retracts by a relative natural
transformation to $C^{\triangleright}_{u_1}$. Hence
$BQ_0$ is contractible. The map $C\to Q_0$ is not
in $W$ if $S^2$ is not $W$-aspherical. But $\alpha$
itself is in $W_\infty\subseteq W$. If it were a
$W$-cofibration, Proposition~\ref{F:couniversal}
would make every cobase change a $W$-equivalence,
contradicting this one. This argument also handles
$W_2$, where the preceding $S^3$ comparison alone
would detect no failure.

For $n=2$ the program computes the actual pushout from the
coend relations. It has eight objects and thirty arrows,
with normalized simplex counts $(8,22,30,20,4)$ and integral
homology ranks $(1,0,0,1,0)$. Every nonzero Smith invariant
of each boundary matrix is $1$. The program checks the
relative deformation functor and natural transformation,
and repeats the calculation for $n=3$. The proof above
establishes all $n\geq2$ without extrapolating those tests.
\end{proof}

\begin{theorem}[Proved: the maximal class allowed by C048]
\label{F:interval-impossibility}
Fix $n\geq2$ and the endpoint inclusion $j$ above. Let $W$
be a fundamental localizer for which $S^2$ is not
$W$-aspherical. If a class $\mathcal K$ of full sieve
inclusions consists of $W$-cofibrations and is closed under
$i\mapsto i\mathbin{\square}j$, then every member of
$\mathcal K$ is an isomorphism. In particular there is no
nontrivial class satisfying C048 for $W_\infty$, for
$W_R$ with a nonzero coefficient ring $R$, or for $W_m$
with $m\geq2$. The class of isomorphisms satisfies the
conditions literally and is the maximal possible class.
\end{theorem}

\begin{proof}
Suppose first that $i:U\hookrightarrow X$ has a crossing
arrow $u\to x$, with $u\in U$ and $x\notin U$. The
classifying functor $\chi:X\to[1]$, equal to $0$ on $U$
and to $1$ on its complement, together with $U\to e$,
is a morphism of arrows from $i$ to $i_0:e\to[1]$.
Choosing $u$ and the indicated crossing arrow gives a
section from $i_0$ to $i$. Their composite is the identity
of $i_0$ in the arrow category. Thus $i_0$ is a retract
of $i$. The pushout-product functor $(-)\mathbin{\square}j$
preserves retract diagrams. If $i\in\mathcal K$, the
retract $i_0\mathbin{\square}j$ of
$i\mathbin{\square}j$ would therefore be a
$W$-cofibration, by Proposition~\ref{F:couniversal}.
This contradicts Theorem~\ref{F:interval-failure}.

If $i$ has no crossing arrows, its sieve property implies
that it is a coproduct-summand inclusion
$U\hookrightarrow U\amalg V$. If it is not an
isomorphism, choose $v\in V$. Its pushout-product with
$j$ is
\[
 U\times I\amalg V\times\{0,2n\}
       \hookrightarrow U\times I\amalg V\times I.
\]
This sieve has the crossing arrow $(v,0)\to(v,1)$.
By closure it belongs to $\mathcal K$, contradicting
the first paragraph. Every member is therefore an
isomorphism. Pushouts and pushout-products of
isomorphisms are isomorphisms, proving the final assertion.
Notice that cobase-change closure was not needed for the
obstruction; adding that requirement cannot repair it.
\end{proof}

This strengthens the calculation at XII ms.~114. The
slice with two incomparable maxima does show failure
of the right-adjoint condition, but the obstruction
already concerns preservation of weak equivalences by
the entire inclusion. There is a related obstruction
to cartesian model structures on $\Cat$ in
\citet[Propositions~4.1--4.3]{F_Pavlov2024}; the theorem
above concerns the precise alternating endpoint
operation in this manuscript. A replacement by a
homotopy pushout, or by a cofibrant replacement of the
ordinary pushout-product map, retains the desired
homotopy type, but changes the specified operation.

\begin{proposition}[Proved: the degree-zero and degree-one exceptions]
\label{F:W0}
For $W_0$, every functor is a $W_0$-cofibration.
For $W_1$, every object-injective functor is a
$W_1$-cofibration. Consequently all full sieve
inclusions form a class satisfying C048 for each
of $W_0$ and $W_1$.
\end{proposition}

\begin{proof}
The functor $\pi_0:\Cat\to\Set$ is left adjoint to
the discrete-category functor and preserves pushouts.
Pushout in sets preserves isomorphisms. Thus pushout
along any functor preserves functors inducing
bijections on components. Full sieves are stable under
cobase change by the collage formula and under the
displayed endpoint pushout-product by inspection of
sources of arrows.

For $W_1$, let $\Pi:\Cat\to\mathbf{Gpd}$ be
groupoid completion, retaining the object set.
It is left adjoint to inclusion and preserves
pushouts. A functor belongs to $W_1$ precisely
when its groupoid completion is an equivalence:
the groupoid $\Pi A$ is the fundamental groupoid
of $BA$ on its vertices. The folk model structure
on categories has equivalences as weak equivalences,
object-injective functors as cofibrations, and is
left proper; these precise properties are recalled
and used in the proof of
\citet[Theorem~1.6]{F_HORR2024}. A pushout of
groupoids in categories is a groupoid, since
its arrows are generated by invertible arrows.
Thus equivalences of groupoids remain equivalences
after cobase change along object-injective
functors. Apply this to the groupoid completions
of the two successive pushouts in the definition
of a $W_1$-cofibration. Object injectivity survives
the first pushout, because object sets of
categorical pushouts are the pushouts of object
sets. It follows that the second comparison
is an equivalence after $\Pi$, as required.
\end{proof}

\begin{theorem}[Proved: equivalent forms of the remaining presheaf
question]\label{F:remaining-equivalence}
The following five assertions, each quantified over all small
categories and all indicated presheaves, are equivalent.
\begin{enumerate}
\item Every $W_\infty$-dwy\'erienne sieve is discretely flat.
\item $W_\infty$-dwy\'erienne sieves are closed under composition.
\item Every totally $W_{\mathbb Z}$-aspherical presheaf is totally
$W_\infty$-aspherical.
\item If $H$ is a totally $W_\infty$-aspherical presheaf on $U$
and $V\subset U$ is a full cosieve with $H|_V$ nonempty, then
$H|_V$ is totally $W_\infty$-aspherical.
\item Every couniversal $W_\infty$-equivalence which is a
full sieve whose complementary full category is isomorphic
to $e$ (one object and only its identity) is
$W_\infty$-dwy\'erienne.
\end{enumerate}
If assertion~2 fails, a counterexample can be chosen with a
classical Dwyer sieve as its first factor and a one-object
extension as its second factor. The equivalence theorem does
not decide whether these five assertions hold.
\end{theorem}

\begin{proof}
We first establish a construction which turns assertion~3
into a precise test. Let $H:V^{\op}\to\Set$ be totally
$W_{\mathbb Z}$-aspherical. Write $E=V/H$; this category is
nonempty, connected, and integrally acyclic. Adjoin an initial
object $0$ to $V$, obtaining $B=V^{\triangleleft}$, and define
$F:B^{\op}\to\Set$ by
\[
 F(0)=\{*\},\qquad F|_V=H.
\]
The restriction maps to $F(0)$ are the unique maps. Its category
of elements is the initial cone $E^{\triangleleft}$, hence is
contractible. We claim that $F$ is totally
$W_\infty$-aspherical.

For a covariant $L:B\to\Set$, put
$D=\operatorname{El}_E((L|_V)\pi)$. The components of $D$ are
integrally acyclic by Proposition~\ref{E:restricted-actions}.
The category
$\operatorname{El}_{B/F}(L\pi_F)$ is obtained from $D$ by
adjoining one source object for each $a\in L(0)$. The arrows
out of that source have crossing element category exactly $E$:
its functor into $D$ sends
\[
 (v,h)\longmapsto (v,h,L(0\to v)(a)).
\]
The uniqueness of $0\to v$ makes this a functor and identifies
its entire crossing category with $E$. Thus the nerve is formed
by attaching one initial cone on $NE$ for each $a$, along these
maps into $ND$. This can either be read degree by degree from
chains starting at a new source or expressed as the corresponding
homotopy pushout. The relative homology of this attachment is
\[
 H_k\bigl(B\operatorname{El}_{B/F}(L\pi_F),BD;\mathbb Z\bigr)
   \cong
 \bigoplus_{a\in L(0)}\widetilde H_{k-1}(BE;\mathbb Z)=0.
\]
Each cone attaches along a nonempty connected category, so it
neither joins distinct components of $D$ nor creates a new one.
If $D$ is empty, nonemptiness of $E$ forces $L(0)$ to be empty
as well. It follows in every case that the resulting components
are acyclic. Proposition~\ref{E:restricted-actions}, together
with contractibility of $B/F$, proves the claim.

Assume assertion~1. Apply it to the sieve obtained by adjoining
one object to $B$ with crossing presheaf $F$. By
Propositions~\ref{F:weight-test} and~\ref{F:DF-actions}, every
restricted action of $B/F$ has contractible components. Take
the covariant indicator of the cosieve $V\subset B$, empty
at $0$ and a singleton on $V$. Its action category is $E$.
Thus $E$ is contractible. The original total
$W_{\mathbb Z}$-asphericity of $H$ supplies all the acyclicity
conditions in Proposition~\ref{E:restricted-actions}; now
$H$ is totally $W_\infty$-aspherical. This proves $1\Rightarrow3$.

For $3\Rightarrow1$, let $H$ be totally $W_\infty$-aspherical
on $U$, and fix $L:U\to\Set$. Let $V=\operatorname{El}_U(L)$
with projection $\rho:V\to U$, and put $H'=\rho^*H$.
There is a canonical isomorphism
\[
 V/H'\cong\operatorname{El}_{U/H}(L\pi)=D.
\]
For any component $D_c$ of $D$, define a subpresheaf
$H'_c\subset H'$ by keeping just those elements whose objects
of $V/H'$ belong to $D_c$. This is a subpresheaf: a restriction
map gives an arrow in the element category, so cannot cross
between components. Its category of elements is exactly $D_c$,
which is nonempty, connected, and acyclic.

For any covariant $M:V\to\Set$, form the covariant $U$-action
\[
 L'(u)=\coprod_{l\in L(u)}M(u,l),
\]
with transitions induced by $L$ and $M$. Then
$\operatorname{El}_{D_c}(M\pi_c)$ is a union of connected
components of $\operatorname{El}_{U/H}(L'\pi)$: both have
objects $(u,h,l,m)$, and restricting to $D_c$ cannot remove
part of a component. The latter components are acyclic by
totality of $H$. Proposition~\ref{E:restricted-actions}
therefore makes $H'_c$ totally $W_{\mathbb Z}$-aspherical.
Assertion~3 makes it totally $W_\infty$-aspherical, so the
identity test makes $D_c$ contractible. This applies to every
$L$ and every component. Proposition~\ref{F:DF-actions}
proves that $H\otimes_U-$ is homotopical; applying this to
every boundary weight proves assertion~1. Empty weights
automatically give the constant empty coend and cause no
exception.

Assertion~3 implies assertion~4 by
Lemma~\ref{F:cosieve-homology}, with $R=\mathbb Z$, followed
by assertion~3 on the restricted presheaf. Assertion~4 implies
assertion~2 by the same restriction and strict-functor
composition argument as in
Theorem~\ref{F:strict-sieve-homology}.

Finally we prove $2\Rightarrow4$ and the stated special form
of a possible counterexample. Given $V\subset U$ and $H$ as
in assertion~4, form $X$ from the old category $U$ and a
disjoint lower copy $V_0$ of $V$ by setting
\[
 X(v_0,u)=U(v,u),\qquad X(u,v_0)=\varnothing
       \quad(v\in V,\ u\in U),
\]
with composition induced by $U$. Since $V$ is a cosieve,
the cosieve generated by $V_0$ in $X$ is the two copies of
$V$ with these crossing arrows, namely $V\times[1]$.
The inclusion $a:V_0\hookrightarrow X$ is classical Dwyer:
the inclusion into this generated cosieve has right adjoint
which forgets the copy. Let $j:U\hookrightarrow X$ be the
old inclusion. The presheaf $G=j_!H$ is total by stability
of totality under Kan extension, and the co-Yoneda formula
gives
\[
 G|_U=H,\qquad G|_{V_0}=H|_V.
\]
Adjoin one object $t$ to $X$, with crossing presheaf $G$,
to obtain a $W_\infty$-dwy\'erienne sieve $b:X\hookrightarrow Y$.
The composite $ba$ has weight $H|_V$ at $t$. Its other
nonempty boundary weights are representables on $V$;
the weights at old objects outside $V$ are empty by the
cosieve condition. Thus $ba$ is $W_\infty$-dwy\'erienne
exactly when $H|_V$ is total. This proves $2\Rightarrow4$.
Applying assertion~4 to $F$ on $V^{\triangleleft}$ constructed
above proves $4\Rightarrow3$.
Proposition~\ref{F:onepoint-couniversal} identifies assertion~5
with assertion~3: a sieve with complement isomorphic to $e$
is exactly the one-point collage of its boundary presheaf. Couniversality
forces that presheaf to be nonempty, and under this condition
the sieve is $W_\infty$-dwy\'erienne precisely when its
presheaf is total $W_\infty$. Thus all five assertions are equivalent, with
the promised restricted form of a failure of composition.
\end{proof}

\begin{proposition}[Open: the remaining class comparisons and
localizers]\label{F:open}
The arguments in this part do not decide
whether every $W_\infty$-dwy\'erienne sieve is
discretely flat, nor whether
$W_\infty$-dwy\'erienne sieves are always closed
under composition. These are equivalent to the
one-point couniversal-recognition question and the
presheaf implication in
Theorem~\ref{F:remaining-equivalence}. Strict and
total functors compose for every fundamental
localizer. For sieves, the universal assertion
of composition is false by
Theorem~\ref{F:low-sieve-composition}; total
composition for a fixed localizer is characterized
exactly by Theorem~\ref{F:total-sieve-cosieves}.
For finite $n\geq2$, strict $W_n$-dwy\'erienne
sieve composition, and total $W_n$-dwy\'erienne
sieve composition without the finite-category
hypothesis, are not settled here. The remaining
$S^2$-aspherical case of C048 is settled in
Corollary~\ref{F:C048-dichotomy} below.
\end{proposition}

The exact missing implication is identified by
Propositions~\ref{E:restricted-actions} and~\ref{F:DF-actions}:
for a total presheaf, must every acyclic component of a
restricted action be simply connected? A covering of such
a component need not itself be an action factoring through
$U/H\to U$, so the covering argument for total categories
does not supply that implication. Equivalently, by
Theorem~\ref{F:remaining-equivalence}, the unresolved question
is whether total integral-homology asphericity of a presheaf
forces its category of elements to be simply connected.
Any obstruction must have a nontrivial perfect fundamental
group. If $U/H$ is totally aspherical,
Proposition~\ref{F:weight-test} already supplies discrete
flatness. The two different interval tests
must be distinguished: the explicit weak-map
comparison has a degree-three obstruction, while
the couniversality test has a degree-two
obstruction and settles $W_2$ by
Proposition~\ref{F:couniversal}. Degree one
behaves differently because groupoid completion
preserves the relevant ordinary pushouts.
